\scanpage{56}
\mspage{62}

En résumé

% margin note written sideways
\margin{\deleted{Si $f_*$ existe, $f$ est une $\unsure{\mathbb D|\underline D}$-équiv. ssi, pour tout faisceau \uuline{constant} $F$ sur $Y$, $F\xrightarrow{\sim} f_*f^*F$ (adj.) est iso. Si $f_!$ existe, $f$ est une $\mathbb D$-équiv. ssi $\forall F$ \uuline{constant} $f_!f^*F\xrightarrow{\sim}F$ (adj) est iso.}}

\uline{Prop. 1} Si \deleted{\unsure{$H^\bullet_{\mathbb D}$}} $p_*$ et $q_*$ existent i.e. si $H^\bullet_{\mathbb D}(X,-)$ et $H^\bullet_{\mathbb D}(Y,\unsure{\cdot|-})$ existent (il suffit $=$ que ça existe pour des coeff. constants) alors $f$ est une $\mathbb D$-équiv. ssi \deleted{\unsure{$\forall N$}} $f$ induit un iso. pour \unsure{la|les} \deleted{\unsure{D}}$\mathbb D$-coh. à coeffts constants de $X$ et $Y$. Si $p_!$ et $q_!$ existent, \inserted{i.e $H_\bullet^{\mathbb D}(X,-)$ et $H_\bullet^{\mathbb D}(Y,-)$ existent} $f$ est une $\mathbb D$-équiv. ssi $f$ induit un iso pour les $\mathbb D$-hom. à coefficients constants de $X$ et $Y$.

\uline{Prop 2} : \deleted{\illegible{1}} les \deleted{\unsure{équivalences}} $\mathbb D$-équiv. Si $X\xrightarrow{f}Y\xrightarrow{g}Z$, si deux parmi $f,g,gf$ sont des $\mathbb D$-équiv, le troisième aussi

\uline{Déf. 2} On dit que $X$ \deleted{\unsure{$\in\mathrm{Ob}$} de $\mathrm{Diag}$} est $\mathbb D$-asphérique (ou $\mathbb D$-contractile, ou $\mathbb D$-triviale) si $X\xrightarrow p e$ est une $\mathbb D$-équivalence, i.e. si $p^*$ est pl. fid.

\deleted{\unsure{Si}} Si $p_*$ existe, cela signifie donc que pour $M\in\mathrm{Ob}\,A$

\[ M\xrightarrow{\overset{\mathrm{adj}}{\sim}}p_*p^*M\qquad\text{i.e. }M\xrightarrow{\sim}H^\bullet_{\mathbb D}(X,M), \]

et si $p_!$ existe, cela signifie que pour tout $M$

\[ p_!p^*M\xrightarrow{\overset{\mathrm{adj}}{\sim}}M\qquad\text{i.e. }H_\bullet^{\mathbb D}(X,M)\xrightarrow{\sim}M. \]

\inserted{Cor.} \deleted{\unsure{NB}} Si $Y$ \deleted{\unsure{$=$}} (resp $X$) est contractile, alors $f:X\longrightarrow Y$ est une $\mathbb D$-équiv. ssi $X$ (resp $Y$) est contractile

\deleted{\uline{\unsure{Déf 3}}}

\uline{Exemple} Prenons $\mathbb D=\mathrm{HOT}$. \inserted{(On peut prendre $\mathrm{Diag}=$ cat des petites catégories)} Alors on a pour $M\in A=\mathrm{Hot}$ et $X$ petite catégorie, \deleted{$p:X\longrightarrow e$}, on a

\[ H^\bullet_{\mathrm{HOT}}(X,M)\simeq\mathrm{Hom}_A(\mathrm{hot}(X),M) \]

où

\[ \mathrm{hot}:\unsure{\{\mathrm{Cat}\}|\mathrm{(Cat)}}\longrightarrow A=\mathrm{Hot} \]

est le foncteur \deleted{qui} can. qui associe à tt $X$ son ``type d'homotopie''. On a

\[ H_\bullet^{\mathrm{HOT}}(X,M)\unsure{=|\simeq}\mathrm{hot}(X)\times M \]

\margin{$\uparrow$ 1} \uuline{Remarques} 1 Pour que $f$ soit $\mathbb D$-équiv, il \uuline{suffit} que pour tt $F$ \uuline{constant} sur $Y$ \inserted{les morphismes d'adjonction} $F\longrightarrow f_*f^*F$ soit iso. \uline{Rem 2} Si $f,f':X\longrightarrow Y$ sont dans une $=$ \unsure{comp.|compos.} connexe de $\uline{\mathrm{Hom}}(X,Y)$, alors $f$ est $\mathbb D$-équiv, $\Leftrightarrow f'$ $\mathbb D$-équiv.
% archive sheet number: 55
